% 行星（综述）
% license CCBYSA3
% type Wiki

本文根据 CC-BY-SA 协议转载翻译自维基百科\href{https://en.wikipedia.org/wiki/Planet}{相关文章}。

\begin{figure}[ht]
\centering
\includegraphics[width=8cm]{./figures/c14ce525d8184ebd.png}
\caption{太阳系八大行星的尺寸比例图（自上而下、从左到右）：土星、木星、天王星、海王星（外行星），地球、金星、火星和水星（内行星）。} \label{fig_Planet_1}
\end{figure}
行星是一种巨大且呈圆形的天体，一般要求其绕恒星、恒星遗迹或褐矮星运行，并且其自身不属于这些天体。按该术语最严格的定义，太阳系共有八颗行星：类地行星（水星、金星、地球和火星）以及巨行星（木星、土星、天王星和海王星）。关于行星形成的最佳理论是星云假说，该假说认为一团星际云从星云中坍缩，形成一颗年轻的原恒星，而这颗原恒星被一个原行星盘所环绕。行星在该盘内通过由引力驱动的物质逐渐累积而成长，这一过程称为吸积。

“行星”（planet）一词源于希腊语 πλανήται（planḗtai），意为“流浪者”。在古代，这个词指的是太阳、月亮，以及在恒星背景上肉眼可见并会移动的五个光点——水星、金星、火星、木星和土星。行星在历史上具有宗教关联：多种文化将天体与神祇对应，而这种与神话和民间传说的联系，也延续到了对新发现的太阳系天体的命名方式之中。随着16至17世纪日心说取代地心说，地球本身也被认作一颗行星。

随着望远镜的发展，“行星”一词的含义被扩展，纳入了只能借助观测工具才能看见的天体：地球之外行星的卫星；冰巨行星天王星与海王星；谷神星及后来被确认属于小行星带的其他天体；以及冥王星，后来被发现是柯伊伯带中冰质天体族群中最大的成员。对柯伊伯带中其他大型天体的发现，尤其是厄里斯，引发了关于如何精确定义“行星”的争论。2006年，国际天文学联合会（IAU）通过了太阳系中“行星”的定义，将四颗类地行星与四颗巨行星归入行星范畴；而谷神星、冥王星与厄里斯则被归入“矮行星”类别。尽管如此，许多行星科学家仍持续以更宽泛的方式使用“行星”一词，包括矮行星以及如月球般呈圆形的卫星。

天文学的进一步发展使得人们发现了超过5900颗太阳系外的行星，即系外行星。这些系外行星常呈现太阳系行星所不具备的特殊性质，例如“热木星”——如飞马座51b般绕其母恒星极近距离运行的巨行星——以及极端偏心的轨道，如HD~20782~b。褐矮星与比木星更大的行星的发现也推动了对于定义边界的讨论，即如何划定行星与恒星之间的分界线。已有多颗系外行星被发现运行在其恒星的宜居带内（即行星表面可能存在液态水的区域），但地球仍是唯一已知能够承载生命的行星。宜居带的形成受多种因素影响，例如行星的大气组成（若其具有大气）。例如火星在赤道区域的白天温度足以形成液态水，只是其气压过低而无法实现。
\subsection{形成}
\begin{figure}[ht]
\centering
\includegraphics[width=8cm]{./figures/edc16278f48809b3.png}
\caption{原行星盘} \label{fig_Planet_2}
\end{figure}
\begin{figure}[ht]
\centering
\includegraphics[width=8cm]{./figures/645312c0cec1ae7e.png}
\caption{行星形成过程中原行星之间的碰撞} \label{fig_Planet_4}
\end{figure}
行星的形成机制目前尚未被完全确定。主流理论认为，行星在星云塌缩形成一层薄的气体和尘埃盘的过程中逐渐聚合而成。在中心区域形成原恒星，其周围则是一个旋转的原行星盘。通过吸积（即通过“粘性碰撞”的过程），盘中的尘埃粒子逐渐聚集质量，形成越来越大的天体。在局部区域中会出现称为微行星的质量集中体，它们通过自身的引力吸引更多物质，从而加速吸积过程。这些质量集中体变得越来越致密，最终在自引力作用下向内塌缩，形成原行星 \(^\text{[6]}\)。

当一个原行星的质量增长到略高于火星的质量时，它会开始吸积一层延展的大气层\(^\text{[7]}\)，从而通过大气阻力效应\(^\text{[8][9]}\)极大提高其对微行星的捕获效率。根据固体与气体的吸积历史不同，最终可能形成巨行星、冰巨行星或类地行星\(^\text{[10][11][12]}\)。人们认为木星、土星和天王星的规则卫星可能以类似方式形成\(^\text{[13][14]}\)；然而，海王星的卫星海卫一很可能是被捕获的\(^\text{[15]}\)。地球的月球\(^\text{[16]}\)与冥王星的卫星卡戎\(^\text{[17]}\)则可能是在巨型碰撞中形成。

当原恒星增大至能够点燃并成为真正的恒星时，残余的吸积盘会从内向外被光致蒸发、太阳风、Poynting–Robertson 阻力及其他效应逐渐清除\(^\text{[18][19]}\)。此后，系统中可能仍存在许多原行星，它们相互环绕或环绕恒星运行，但随着时间推移，其中许多会发生碰撞，或合并成更大的原行星，或为其他原行星提供可吸积的碎片\(^\text{[20]}\)。那些质量足够大的天体会吸尽其轨道邻域的大部分物质而成为行星。避免了碰撞的原行星可能通过引力捕获成为行星的天然卫星，或者留存在某些轨道带中，最终成为矮行星或小天体\(^\text{[21][22]}\)。
\begin{figure}[ht]
\centering
\includegraphics[width=8cm]{./figures/624858ff17708327.png}
\caption{超新星遗迹喷 ejecta 产生行星形成物质} \label{fig_Planet_5}
\end{figure}
较小的微行星的高能撞击（以及放射性衰变）会加热正在成长的行星，使其至少部分熔化。行星内部开始按照密度分异，密度更高的物质下沉至核心。\(^\text{[23]}\)较小的类地行星在这种聚积过程中会失去其大部分大气，但失去的气体可通过地幔逸气以及随后彗星的撞击所补充\(^\text{[24]}\)（更小的行星会通过各种逃逸机制而失去它们获得的任何大气\(^\text{[25]}\)）。

随着对太阳以外恒星的行星系统的发现与观测，人们得以对上述理论加以扩展、修正甚至替代。金属丰度——一个天文学术语，用来描述原子序数大于 2（即氦）的化学元素的丰度——似乎决定着恒星拥有行星的可能性。\(^\text{[26][27]}\)因此，与金属贫乏的第二星族（Population II）恒星相比，金属丰富的第一星族（Population I）恒星更可能拥有规模可观的行星系统。\(^\text{[28]}\)
\subsection{太阳系中的行星}
根据 IAU 的定义，太阳系中共有八颗行星（按距太阳距离递增排列）\(^\text{[2]}\)：Mercury、Venus、Earth、Mars、Jupiter、Saturn、Uranus 与 Neptune。Jupiter 质量最大，约为地球质量的 318 倍；而 Mercury 最小，仅为地球质量的 0.055 倍。\(^\text{[29]}\)

太阳系行星可依据组成成分划分为不同类别。类地行星（terrestrials）与 Earth 类似，其主体由岩石与金属构成：Mercury、Venus、Earth 与 Mars。Earth 是类地行星中体积最大的一个。\(^\text{[30]}\)

巨行星（giant planets）的质量显著大于类地行星：Jupiter、Saturn、Uranus 与 Neptune。\(^\text{[30]}\)它们在成分上与类地行星存在根本差异。气体巨行星（gas giants）Jupiter 与 Saturn 主要由氢与氦组成，是太阳系中最为巨大的行星。Saturn 的质量约为 Jupiter 的三分之一，即约 95 个地球质量。\(^\text{[31]}\)

冰巨行星（ice giants）Uranus 与 Neptune 则主要由低沸点物质组成，如 water、methane 与 ammonia，并拥有由 hydrogen 与 helium 构成的厚大气层。它们的质量显著低于气体巨行星，仅分别为 14 与 17 个地球质量。\(^\text{[31]}\)
\begin{figure}[ht]
\centering
\includegraphics[width=8cm]{./figures/f984108e351cc1ba.png}
\caption{太阳、行星、矮行星与卫星的尺寸按比例标注示意。各天体之间的距离不按比例绘制。小行星带位于火星与木星轨道之间，柯伊伯带位于海王星轨道之外。} \label{fig_Planet_6}
\end{figure}
\noindent 矮行星是达到引力塑形（即具有趋于球形的外形），但尚未清除其轨道附近其他天体的天体。按照距太阳平均距离递增排列，被天文学界普遍认可的矮行星包括：Ceres，Orcus，Pluto，Haumea，Quaoar，Makemake，Gonggong，Eris 与 Sedna。\(^\text{[32][33]}\)Ceres 位于 Mars 与 Jupiter 轨道之间的小行星带，是该区域最大的天体。其余八个均位于 Neptune 轨道之外。Orcus、Pluto、Haumea、Quaoar 与 Makemake 运行于 Kuiper 带——这是位于 Neptune 轨道之外的第二个太阳系小天体带。Gonggong 与 Eris 运行于散射盘（scattered disc），其位置更为外侧，且该区域不同于 Kuiper 带——会因 Neptune 的引力扰动而呈现轨道不稳定性。Sedna 是目前已知最大的游离天体（detached object），其轨道不会接近太阳到足以受任何经典行星的显著影响；其轨道成因仍在讨论之中。以上九个天体均与类地行星相似，拥有固体表面，但其成分主要是冰与岩石，而非岩石与金属。此外，它们全部小于 Mercury，其中 Pluto 是目前已知体积最大的矮行星，而 Eris 的质量最大。\(^\text{[34][35]}\)

至少有十九颗具行星质量的卫星（planetary-mass moons 或 satellite planets），即质量足以呈现椭球形状的卫星\(^\text{[4]}\)：

\begin{itemize}
    \item Earth 的一颗卫星：Moon；
    \item Jupiter 的四颗卫星：Io，Europa，Ganymede 与 Callisto；
    \item Saturn 的七颗卫星：Mimas，Enceladus，Tethys，Dione，Rhea，Titan 与 Iapetus；
    \item Uranus 的五颗卫星：Miranda，Ariel，Umbriel，Titania 与 Oberon；
    \item Neptune 的一颗卫星：Triton；
    \item Pluto 的一颗卫星：Charon。
\end{itemize}
Moon、Io 与 Europa 的成分与类地行星相似；其它卫星的成分则与矮行星类似，以冰与岩石为主，其中 Tethys 几乎完全由冰组成。由于其表面的冰层使内部难以观测，Europa 常被视为一颗冰质行星。\(^\text{[4][36]}\)

Ganymede 与 Titan 在半径上均大于 Mercury，而 Callisto 也几乎与之相当，但这三者的质量均远小于 Mercury。Mimas 是公认的几何意义上的行星（即达到地质行星标准）中最小的一个，其质量约为地球质量的六百万分之一；然而，还有许多比它更大的天体却可能并非地质行星（例如 Salacia）。\(^\text{[32]}\)
\subsection{Exoplanets}
\begin{figure}[ht]
\centering
\includegraphics[width=8cm]{./figures/85a4bfec6f369e50.png}
\caption{截至 2023 年 8 月的每年系外行星探测数量（数据来自 NASA 系外行星档案）\(^\text{[37]}\)} \label{fig_Planet_8}
\end{figure}
系外行星（exoplanet）指位于太阳系之外的行星。截至 2025 年 10 月 30 日，已经确认的系外行星共有 6{,}128 颗，分布在 4{,}584 个行星系统中，其中 1{,}017 个系统拥有不止一颗行星。\(^\text{[38]}\)已知的系外行星尺寸跨度极大，从比 Jupiter 大两倍的气态巨行星到仅略大于 Moon 的微小行星均有分布。对引力微透镜（gravitational microlensing）数据的分析表明，在 Milky Way 中，每颗恒星平均至少拥有 1.6 颗束缚行星。\(^\text{[39]}\)

1992 年初，无线电天文学家 Aleksander Wolszczan 与 Dale Frail 宣布发现两颗绕脉冲星 PSR 1257+12 运行的行星。\(^\text{[40]}\)这项发现随后得到证实，并普遍被认为是首次确定无疑的系外行星探测。研究者推测，这些行星可能形成于产生该脉冲星的超新星爆炸后残留的盘状物质中。\(^\text{[41]}\)

首次确认的、绕普通主序星运行的系外行星于 1995 年 10 月 6 日被发现。当时 Geneva 大学的 Michel Mayor 与 Didier Queloz 宣布在恒星 51 Pegasi 周围探测到行星 51 Pegasi b。\(^\text{[42]}\)从那时直到 Kepler 空间望远镜任务开始前，人类已知的大多数系外行星都是质量与 Jupiter 相当或更大的气态巨行星，因为它们更易被探测到。Kepler 候选行星目录中大多数行星的尺度介于 Neptune 与更小的行星之间，有些甚至小于 Mercury。\(^\text{[43][44]}\)

2011 年，Kepler 团队报告了首次发现绕类太阳恒星运行的地球大小行星：Kepler-20e 与 Kepler-20f。\(^\text{[45][46][47]}\)自那时起，已有超过 100 颗大小近似 Earth 的行星被确认，其中有 20 颗位于其恒星的宜居带（habitable zone）内——这是允许类地行星在具备足够大气压的条件下，其表面可能存在液态水的轨道范围。\(^\text{[48][49][50]}\)据估计，每五颗类太阳恒星中就有一颗拥有位于宜居带中的地球大小行星，这意味着距离地球最近的此类行星预计不会超过 12 光年。\(^\text{[a]}\)此类类地行星的出现频率是 Drake 方程中的变量之一，该方程用于估算 Milky Way 中存在的具备智慧并能够通信的文明数量。\(^\text{[53]}\)

在太阳系中不存在的行星类型包括超地球（super-Earth）与迷你海王星（mini-Neptune），其质量介于 Earth 与 Neptune 之间。质量不超过约两倍地球质量的天体通常被认为具有类似 Earth 的岩石组成；超过该范围，则会演化为含有大量挥发物与气体、类似 Neptune 的结构。\(^\text{[54]}\)行星 Gliese~581c 的质量为 Earth 的 $5.5$--$10.4$ 倍\(^\text{[55]}\)，在被发现时因可能处于其恒星的宜居带而引起关注\(^\text{[56]}\)，但后续研究表明它实际上距离其恒星过近，无法宜居。\(^\text{[57]}\)更大质量的行星亦已被发现，其性质与质量向上连续延伸至褐矮星（brown dwarf）领域。\(^\text{[58]}\)

已知一些系外行星与其母恒星之间的距离远小于太阳系中任意一颗行星与 Sun 的距离。太阳系内距离 Sun 最近的 Mercury 距离为 $0.4\ \mathrm{AU}$，其公转周期为 $88$ 天；但超短周期行星（ultra-short-period planets）可以在不到一天内完成公转。Kepler-11 系统中有五颗行星的轨道周期都短于 Mercury，而且它们的质量都远大于 Mercury。  
另有热木星（hot Jupiters），例如 51~Pegasi~b\(^\text{[42]}\)，紧贴恒星运行，可能因蒸发而变成“冥狱行星”（chthonian planets），也就是仅剩下残核的行星。  
此外，还有系外行星比太阳系中任何行星离其恒星更远。Neptune 距 Sun 为 $30$ AU，公转周期为 $165$ 年；但某些系外行星距离其恒星数千 AU，公转周期超过百万年（如 COCONUTS-2b）。\(^\text{[59]}\)
\subsection{属性}
尽管每颗行星都具有各自独特的物理特征，但它们之间仍然存在一些普遍的共同点。其中一些特征——如行星环或天然卫星——目前仅在太阳系行星中被观测到，而另一些则在系外行星中也普遍存在。\(^\text{[60]}\)
\subsection{动力学特征}
\textbf{轨道（Orbit）}
\begin{figure}[ht]
\centering
\includegraphics[width=8cm]{./figures/0d173cd711ffbba6.png}
\caption{行星海王星的轨道与冥王星轨道的比较。请注意冥王星轨道相对于海王星轨道的拉长（偏心率），以及其相对于黄道面的巨大夹角（轨道倾角）。} \label{fig_Planet_7}
\end{figure}
在太阳系中，所有行星绕日公转的方向都与太阳自转方向一致：从太阳北极上方看为逆时针方向。至少已经发现一颗系外行星（WASP-17b）其公转方向与其恒星自转方向相反。\(^\text{[61]}\)一个行星绕恒星运行一周的周期称为其恒星周期或“年”。\(^\text{[62]}\)行星的“年长”取决于它与恒星的距离：距离越远，轨道越长，同时由于恒星引力作用较弱，其运行速度也越慢。

没有任何行星的轨道是完全圆形的，因此其与恒星的距离会在一年内发生变化。行星最接近恒星的位置称为“近日点”（在太阳系中称“近日点”）；最远的位置称为“远日点”。当行星接近近日点时，其速度会增大，因为它将引力势能转化为动能，就像地球上自由落体物体下落时会加速一样。当地行星接近远日点时，其速度会减小，就像向上抛出的物体在接近最高点时速度会逐渐减小一样。\(^\text{[63]}\)

每个行星的轨道由一组轨道要素来描述：
\begin{itemize}
    \item 偏心率（eccentricity）描述行星椭圆（椭圆形）轨道的扁长程度。偏心率低的行星轨道更接近圆形，而偏心率高的行星轨道则更为椭圆。太阳系中的行星及大型卫星的偏心率都相对较低，因此它们的轨道几乎是圆形的。\(^\text{[62]}\)彗星、许多柯伊伯带天体以及若干系外行星则具有非常高的偏心率，因此其轨道极为椭圆。\(^\text{[64][65]}\)
    \item 半长轴（semi-major axis）给出轨道的尺度。它是椭圆轨道最长直径的一半。此值并不等同于远日点距离，因为没有任何行星的轨道是以恒星为椭圆几何中心的。\(^\text{[62]}\)
    \item 轨道倾角（inclination）表示行星轨道相对于某一参考平面的倾斜程度。在太阳系中，参考平面是地球轨道平面，称为黄道面。对于系外行星，参考平面称为“天空平面”，即垂直于地球观察方向的平面。\(^\text{[66]}\)太阳系八大行星的轨道都非常接近黄道面；然而，一些较小的天体，如智神星（Pallas）、冥王星（Pluto）和阋神星（Eris），以及彗星的轨道倾角都要大得多。\(^\text{[67]}\)大型卫星的轨道通常也不太偏离其母行星的赤道，但地球的月球、土星的土卫八（Iapetus）以及海王星的海卫一（Triton）是例外。其中海卫一尤为独特，它是唯一逆行绕行母行星的大型卫星，即其轨道方向与母行星自转方向相反。\(^\text{[68]}\)
    \item 行星轨道穿过参考平面的两个点称为升交点（ascending node）和降交点（descending node）。\(^\text{[62]}\)升交点经度（longitude of the ascending node）是从参考平面零经度到升交点的角度。近日点参数（argument of periapsis，太阳系中称近日点参数）是从升交点到行星最接近恒星位置之间的角度。\(^\text{[62]}\)
\end{itemize}
\textbf{自转轴倾角}
\begin{figure}[ht]
\centering
\includegraphics[width=8cm]{./figures/9e2bb5a76a5b4400.png}
\caption{地球的自转轴倾角约为 $23.4^\circ$。它在 $22.1^\circ$ 与 $24.5^\circ$ 之间以约 $41{,}000$ 年为周期振荡，目前正处于减小阶段。} \label{fig_Planet_9}
\end{figure}
行星具有不同程度的自转轴倾角；它们相对于其恒星赤道平面呈一定角度自转。这导致行星表面各半球在其公转周期内所接收的光照量发生变化：当北半球背向恒星时，南半球则朝向恒星，反之亦然。因此，每颗行星都会经历季节，其气候会随着一年中的位置变化而改变。每个半球在其自转轴最远离或最近指向恒星的位置，称为至点。每颗行星在一个公转周期中都有两个至点；当一个半球处于夏至且白昼时间最长时，另一半球则处于冬至且白昼最短。各半球所接收的光和热量的差异，会在每年引起对应半球天气模式的变化。

木星的自转轴倾角非常小，因此其季节变化极其微弱；而天王星的倾角则极端到几乎“侧躺”着自转，这意味着其半球在至点附近要么持续处于日照，要么持续处于黑夜。\(^\text{[69]}\)在太阳系中，水星、金星、谷神星和木星的倾角非常小；智神星、天王星和冥王星的倾角极大；而地球、火星、灶神星、土星和海王星的倾角属于中等范围。\(^\text{[70][71][72][73]}\)

对于系外行星，其自转轴倾角目前尚不确定，但由于靠近母恒星，大多数“热木星”被认为倾角极小。\(^\text{[74]}\)类似地，行星质量的卫星自转轴倾角通常接近零，\(^\text{[75]}\)其中地球的月球以 $6.687^\circ$ 为最大例外；\(^\text{[76]}\)此外，木卫四（卡利斯托）的自转轴倾角在数千年的时间尺度上会在 $0^\circ$ 到约 $2^\circ$ 之间变化。\(^\text{[77]}\)

\textbf{自转}

行星围绕穿过其中心的不可见自转轴进行自转。行星自转一周的时间称为恒星日（stellar~day）。太阳系中大多数行星的自转方向与它们绕太阳公转的方向一致，即从太阳北极上方观察为逆时针方向。然而，金星\(^\text{[78]}\)与天王星\(^\text{[79]}\)是例外，它们呈顺时针方向自转。不过，由于天王星自转轴倾角极端，对于其“北极”方向存在不同的约定，因此关于它是顺时针还是逆时针自转也存在不同定义。\(^\text{[80]}\)无论采用哪种约定，天王星相对于其轨道而言都属于逆行自转（retrograde~rotation）。\(^\text{[79]}\)
\begin{figure}[ht]
\centering
\includegraphics[width=8cm]{./figures/0c97e1c28cdd96c6.png}
\caption{行星与月球的自转周期（加速 10\,000 倍；负值表示逆行自转）、扁率以及自转轴倾角的比较（SVG 动画）} \label{fig_Planet_10}
\end{figure}
行星的自转可在形成过程中由多种因素诱发。吸积物体各自所携带的角动量会汇总为行星的总体角动量；巨行星吸积气体的过程也会增加其自转角动量；最后，在行星形成后期，原行星之间的随机吸积碰撞会随机改变行星的自转轴方向。\(^\text{[81]}\)

行星昼长差异极大：金星完成一次自转需要 243 天，而巨行星仅需数小时。\(^\text{[82]}\)系外行星的自转周期通常未知，但对“热木星”而言，由于其极度接近母恒星，它们会被潮汐锁定（即其自转与公转同步）。因此，它们始终以同一面朝向恒星，一侧永昼，一侧永夜。\(^\text{[83]}\) 水星与金星（最靠近日面的两颗行星）同样具有非常缓慢的自转：水星被潮汐锁定在 3:2 的自旋–轨道共振中（每绕太阳两周自转三次），\(^\text{[84]}\)而金星的自转可能处于潮汐力减速与由太阳加热驱动的大气潮汐加速之间的一种平衡状态。\(^\text{[85][86]}\)

所有的大型卫星都被潮汐锁定于其母行星；\(^\text{[87]}\)冥王星与其卫星卡戎互相潮汐锁定，\(^\text{[88]}\)妖神星（Eris）与其卫星迪斯诺米亚（Dysnomia）亦然，\(^\text{[89]}\)甚至可能包括阎王星（Orcus）与其卫星 Vanth。\(^\text{[90]}\)其他已知自转周期的矮行星自转速度都比地球快；妊神星（Haumea）的自转极快，以至被拉伸成一个三轴椭球体。\(^\text{[91]}\)系外行星 Tau Boötis b 与其母星 Tau Boötis 似乎也互相潮汐锁定。\(^\text{[92][93]}\)

\textbf{轨道清空}

根据国际天文学联合会（IAU）的定义，行星的关键动力学特征在于其已经清空了自身轨道邻域。一颗已清空邻域的行星，其质量足够大，能够聚集或清扫掉位于其轨道上的所有微行星（planetesimals）。实际上，它是独立围绕恒星运行的，而不是与大量尺寸相近的天体共享同一轨道。如上所述，这一特征被纳入 IAU 于 2006 年 8 月通过的行星正式定义中。[2] 虽然这一标准目前仅适用于太阳系，但已经发现若干年轻的系外行星系统，其原恒星盘（circumstellar disc）中存在能量或结构迹象，暗示轨道清空过程正在发生。[94]
\subsubsection{物理特征}
\textbf{大小与形状}

行星在自身重力作用下被拉向近似球形，因此行星的大小可以用平均半径来粗略表达（例如地球半径、木星半径）。然而，行星并非完美球体；例如，由于地球自转，其两极略微扁平，赤道附近略有鼓起。[95] 因此，用扁球体（oblate spheroid）来近似地球形状更为准确，其赤道直径比极轴直径大约 43 公里（27 英里）。[96]

一般而言，行星的形状可以通过给出扁球体的极半径与赤道半径，或指定一个参考椭球体来描述。由此可以计算行星的扁率、表面积和体积；进一步结合其大小、形状、自转速率以及质量，即可计算其正常重力。[97]

\textbf{质量}

行星最重要的物理特征在于：其质量足够大，使自身重力能够压倒维系其物理结构的电磁力，使其达到\textbf{流体静力学平衡（hydrostatic equilibrium）}。这意味着所有行星都是球形或球状体。低于某一质量上限时，天体可以保持不规则形状；但当其化学组成所允许的临界质量被超过时，重力会开始将物体拉向自身质心，使其最终塌缩成球体。[98]

质量也是区分行星与恒星的主要属性。太阳与木星之间的质量区间在太阳系中无任何天体，但系外行星中存在此类天体。恒星的最低质量极限约为 75–80 个木星质量（MJ）。部分学者提议把这一界限作为行星的最大质量上限，因为从“显著自压缩开始”（约一土星质量）一直到氢点火成为红矮星，其内部物理并无本质变化。[54]另一些观点认为约 13 MJ（太阳丰度条件下）可以作为行星与褐矮星的分界，因为在此质量上天体能够发生氘核聚变。[99] 然而氘燃烧持续时间短，大多数褐矮星早已结束其氘燃烧阶段。[58] 这一界限并未得到普遍认可：Exoplanets Encyclopaedia 将上限取到 60 MJ，Exoplanet Data Explorer 则取到 24 MJ。[100][101]

已知质量最小、质量测量可靠的系外行星是 1992 年发现的脉冲星行星 PSR B1257+12A，其质量约为水星的一半。[102] 更小的是围绕白矮星运行的 WD 1145+017 b，其质量约与矮行星妊神星（Haumea）相当，通常被称为小行星（minor planet）。[103]围绕主序星（非太阳）运转、质量最小的已知行星是 Kepler-37b，其质量（和半径）可能略大于月球。[44] 太阳系中普遍认可的最小地球物理行星是土星的卫星米玛斯（Mimas），半径约为地球的 3.1\%，质量约为地球的 0.00063\%。[104] 土星的另一颗更小的卫星——福柏（Phoebe），目前是不规则天体，半径为地球的 1.7\%，质量为 0.00014\%，但据信其在早期曾达到流体静力平衡并发生内部分异，之后被撞击破坏而变形。[105][104][106]部分小行星可能是开始聚积并发生分异的原行星（protoplanets）的碎块，但因灾难性碰撞，仅存金属核或岩质核心；[107][108][109] 或者是这些碰撞碎片再聚集形成的天体。[110]

\textbf{内部分异}
\begin{figure}[ht]
\centering
\includegraphics[width=8cm]{./figures/a7021b87c5306863.png}
\caption{木星内部结构示意图：在岩石核心之上覆盖着深厚的金属氢层} \label{fig_Planet_11}
\end{figure}
每一颗行星在起初形成时都处于完全流体状态；在形成早期，更致密、更重的物质下沉至中心，使较轻的物质留在靠近表面的区域。因此，每颗行星都拥有分异的内部结构，由致密的行星核心组成，外围包围着一个曾经或仍然处于流体状态的地幔。类地行星的地幔被坚硬的地壳封闭在内部[111]，但在巨行星中，地幔与其上方的云层逐渐过渡并融合。类地行星的核心主要由铁和镍等元素构成，其地幔则主要由硅酸盐组成。木星和土星被认为拥有由岩石与金属构成的核心，外包一层金属氢地幔[112]。天王星和海王星体积较小，它们拥有岩石核心，外围包围着由水、氨、甲烷及其他冰类物质构成的地幔[113]。这些行星核心内部的流体运动产生了地磁发电机效应，从而生成行星的磁场[111]。

类似的内部分异过程也被认为发生在部分大型卫星与矮行星上[32]，但这一过程并非总能完全进行：谷神星（Ceres）、卡利斯托（Callisto）以及泰坦（Titan）似乎都仅完成了部分分异[114][115]。小行星灶神星（Vesta）虽然因受到撞击破坏而失去球形、未被归类为矮行星，但其内部结构已经完成分异[116]，与金星、地球和火星的情况相似[109]。

\textbf{大气}
\begin{figure}[ht]
\centering
\includegraphics[width=8cm]{./figures/01c321ee2d5369eb.png}
\caption{“地球大气层”} \label{fig_Planet_12}
\end{figure}
太阳系中除水星以外的所有行星都拥有相当稠密的大气层[117]，因为它们的引力足够强，可以将气体束缚在靠近表面的区域。土星最大的卫星——泰坦（Titan）也拥有稠密的大气层，其厚度甚至超过地球大气层[118]；海王星最大的卫星——海卫一（Triton）[119]以及矮行星冥王星则拥有更为稀薄的大气层[120]。体积更大的巨行星质量足以维持大量轻质气体（如氢和氦），而体积较小的行星则会将这些气体流失至太空[121]。对系外行星的分析表明，保留这些轻质气体的临界质量约为 $2.0^{+0.7}_{-0.6}\,M_{\oplus}$，因此地球与金星大致位于岩质行星大小的上限附近[54]。

地球大气的成分与其他行星不同，因为生命过程向大气中引入了游离的分子氧[122]。火星与金星的大气都以二氧化碳为主，但其密度差异极大：火星平均表面气压不到地球的 1\%（不足以允许液态水存在）[123]，而金星平均表面气压约为地球的 92 倍[124]。金星目前极端的高温高压环境被认为源于其历史上的失控温室效应，使其成为太阳系中表面温度最高的行星，甚至超过更靠近太阳的水星[125]。尽管金星表面条件极端严酷，其大气中约 50–55 公里高度的温度与压力却与地球相近（这是太阳系中除地球外唯一存在此类条件的区域），因而该高度区域被提出为未来人类探索的可能基地[126]。泰坦拥有除地球外太阳系中唯一富含氮气的稠密行星大气。类似于地球表面条件接近水的三相点，使水能在三种物态间循环；泰坦的环境条件则接近甲烷的三相点[127]。

行星大气在不同程度上受入射辐射或内部能量驱动，导致其形成动态的天气系统，例如：地球上的飓风、火星上的全球性沙尘暴、木星上比地球还大的反气旋（大红斑），以及海王星大气中的空洞[69]。在系外行星上观察到的天气现象包括：在 HD 189733 b 上出现的、面积为木星大红斑两倍的高温区域[128]；以及位于热木星 Kepler-7b、超级地球 Gliese 1214 b 等行星上的云层[129][130][131]。

由于热木星与其宿主恒星极端接近，它们的大气在恒星辐射下不断被剥离，如同彗星般形成尾迹[132][133]。这类行星常在昼夜两侧呈现巨大温差，从而可能产生超音速风流[134]；然而，影响昼夜温差与大气动力学的因素众多，其细节十分复杂[135][136]。

\textbf{磁层}
\begin{figure}[ht]
\centering
\includegraphics[width=8cm]{./figures/19b83eb33522eada.png}
\caption{地球磁层（示意图）} \label{fig_Planet_13}
\end{figure}
行星的一项重要特征是其本征磁矩，而本征磁矩又会产生磁层。磁场的存在表明行星在地质学意义上仍然“活着”。换言之，具有磁化性质的行星在其内部存在电导流体的流动，从而产生磁场。这些磁场会显著改变行星与太阳风之间的相互作用。一个具有磁化的行星会在其周围的太阳风中开辟出一个空腔，这个区域称为磁层（magnetosphere），太阳风无法穿透其中。磁层的尺寸可能远大于行星本身。相反，非磁化行星只能依靠电离层与太阳风的相互作用而形成较小的诱导磁层，这种磁层无法有效保护行星[137]。

在太阳系八大行星中，只有金星和火星缺乏这样的磁场[137]。在具有磁场的行星中，水星的磁场最弱，几乎无法偏转太阳风。木星的卫星伽利略卫星之一——盖尼米德（Ganymede）拥有更强的磁场，而木星自身的磁场是太阳系最强的（其强度之大，事实上对未来所有前往其轨道内卡利斯托以内的载人任务构成严重的健康风险[138]）。其他巨行星在其表面的磁场强度与地球相近，但它们的磁矩却显著更大。天王星和海王星的磁场相对于其自转轴具有强烈的倾角，并且磁场中心相对于行星中心存在明显偏移[137]。

2003年，夏威夷的一组天文学家在观测恒星 HD 179949 时发现其表面出现一个亮斑，显然是由其轨道上一颗热木星的磁层所诱发的[139][140]。
\subsubsection{次级特征}
\begin{figure}[ht]
\centering
\includegraphics[width=8cm]{./figures/7eabc59da7d0513c.png}
\caption{土星的行星环} \label{fig_Planet_15}
\end{figure}
在太阳系中，若干行星或矮行星（例如海王星和冥王星）具有彼此之间，或与更小天体之间的轨道周期共振现象。这种情况在卫星系统中尤为常见（例如木星的木卫一、木卫二与木卫三之间的共振，或土星的土卫二与土卫四之间的共振）。除水星与金星之外，所有行星都拥有天然卫星（通常称为“月亮”）。地球有一颗卫星，火星有两颗，而巨行星拥有数量众多、结构复杂的“类行星系统”般的卫星群。除谷神星（Ceres）与赛德娜（Sedna）外，所有共识矮行星也都至少拥有一颗卫星。许多巨行星的卫星与类地行星及矮行星具有相似的地质特征，其中一些被研究为可能适宜生命存在的地点（尤其是木卫二与土卫二）[141][142][143][144][145]。

四大巨行星都拥有不同规模和复杂度的行星环。这些环主要由尘埃或颗粒组成，但可能包含微小的“卫星状微体”（moonlets），其引力有助于塑造并维持环的结构。虽然行星环的确切起源尚未明晰，但普遍认为它们源自掉入母行星洛希极限（Roche limit）内部、因潮汐力而被撕裂的天然卫星[146][147]。矮行星妊神星（Haumea）[148]和夸欧尔（Quaoar）[149]也被发现拥有行星环。

在系外行星周围尚未观测到类似的次级结构。亚褐矮星 Cha 110913−773444（常被描述为“流浪行星”）被认为拥有一个微小的原行星盘\,[150]，而亚褐矮星 OTS 44 被证实拥有质量至少相当于十个地球质量的巨大原行星盘\,[151]。
\subsection{历史与词源}
行星的概念在天文学历史上不断演变，从古代被视为“神圣之光”的天体，发展为科学时代的类地天体。该概念也已从仅指太阳系内的世界扩展至无数系外行星系统。关于“哪些天体应当被视为行星”这一共识曾多次改变：历史上，行星的范畴曾包括小行星、卫星以及像冥王星这样的矮行星[152][153][154]，而时至今日，这一问题仍未完全达成一致[154]。
\subsubsection{古代文明与古典行星}
由于太阳系中的五颗古典行星可被肉眼直接观测到，它们自古以来便为人类所知，并对神话、宗教宇宙观以及古代天文学产生了深远影响。古代天文学家注意到，某些光点会在天空中移动，与保持相对位置恒定的“恒星”不同。[155] 古希腊人将这些移动的光点称为 πλάνητες ἀστέρες（planētes asteres，“游移之星”），或简称为 πλανῆται（planētai，“游者”），[156] 今日的术语“planet（行星）”便源自此。[157][158][159]在古希腊、中国、巴比伦，以及几乎所有前现代文明中，[160][161] 人们普遍相信地球是宇宙的中心，所有的“行星”都绕地球运行。形成这一认知的原因包括：恒星与行星似乎每天都绕着地球旋转，[162] 以及基于常识的直观判断——地球坚实而稳定，没有运动迹象，似乎处于静止状态。[163]

\textbf{巴比伦}

已知最早拥有成熟行星理论的文明是巴比伦人，他们在公元前第一与第二千纪居住于美索不达米亚。现存最古老的行星天文文本为《阿米萨杜卡的金星泥板》（\emph{Babylonian Venus tablet of Ammisaduqa}），这是一份公元前7世纪的抄本，记录了对金星运动的观测，其原本可能早至公元前第二千纪。[164]《MUL.APIN》是两块源自公元前7世纪的楔形文字泥板，系统展示了太阳、月亮与行星在一年中的运行规律。[165]巴比伦后期的天文学是西方天文学乃至整个西方精确科学传统的起源。[166]《Enuma anu enlil》，写于新亚述时期（公元前7世纪），[167] 包含了一份凶兆列表以及它们与各种天象（包括行星运动）的对应关系。[168][169] 巴比伦天文学家识别了内行星金星、 水星，以及外行星火星、木星、土星，这些在望远镜发明前一直是人类已知的全部行星。[170]

\textbf{希腊—罗马天文学}

古希腊人最初并不像巴比伦人那样强调行星的重要性。公元前6至5世纪，毕达哥拉斯学派似乎发展了独立的行星体系，其中地球、太阳、月亮与行星都围绕宇宙中心的“中央之火”（\emph{Central Fire}）运行。据说毕达哥拉斯或巴门尼德斯最早意识到昏星（\emph{Hesperos}）与晨星（\emph{Phosphoros}）实际上是同一颗星（阿佛洛狄忒，即后来的金星），[171] 尽管这一点在美索不达米亚早已为人所知。[172][173] 公元前3世纪，萨摩斯的阿里斯塔克斯提出了日心体系，认为地球与行星绕太阳运行。然而地心体系一直占主导地位直至科学革命时期。[163]

到公元前1世纪的希腊化时期，希腊人开始发展自己的数学方法来预测行星位置。这些基于几何学而非巴比伦人算术方法的理论，最终在复杂度与完备性上超越了巴比伦体系，能够解释肉眼可见的绝大部分天体运动。该体系的最高成就在公元2世纪托勒密的《天文学大成》（\emph{Almagest}）中实现，其影响力巨大，以至于取代了所有更早的天文学著作，并在西方世界作为权威天文文本持续了13个世纪。[164][174]对希腊人与罗马人而言，可见的“行星”共有七个，它们都被认为依照托勒密定义的复杂运动规律绕地球运行。按托勒密的顺序并使用现代名称，这七个天体依次是：月球、水星、金星、太阳、火星、木星与土星。[159][174][175]
\subsubsection{中世纪天文学}
\begin{figure}[ht]
\centering
\includegraphics[width=8cm]{./figures/200a04b44cceda48.png}
\caption{1660 年克劳狄乌斯·托勒密地心模型的插图} \label{fig_Planet_16}
\end{figure}
在西罗马帝国灭亡之后，天文学在印度和中世纪伊斯兰世界中继续发展。公元 499 年，印度天文学家阿耶波多（Aryabhata）提出了一个行星模型，其中明确包含地球绕自身轴心的自转，并将这种自转解释为恒星看似向西移动的原因。他还推测行星运行轨道为椭圆形 [176]。阿耶波多的追随者在印度南部尤为兴盛，他们遵循他的地球自转等原理，并基于这些思想撰写了大量后续著作 [177]。

伊斯兰黄金时代的天文学主要在中东、中亚、安达卢斯与北非开展，之后传播至远东与印度。这些天文学家（如多才多艺的伊本·海赛姆 Ibn~al-Haytham）通常接受地心宇宙模型，但他们质疑托勒密的本轮体系并试图提出替代理论。10 世纪的天文学家阿布·赛义德·西吉（Abu~Sa'id~al-Sijzi）接受了地球绕自转轴自转的观点 [178]。11 世纪时，伊本·西那（Avicenna）观测到了金星凌日现象 [179]。其同时代学者比鲁尼（Al-Biruni）利用三角学提出了一种测量地球半径的新方法，与埃拉托色尼的方法不同，这一方法只需在一座高山进行观测即可 [180]。
\subsubsection{科学革命与外行星的发现}
\begin{figure}[ht]
\centering
\includegraphics[width=8cm]{./figures/43a174f81fae1e03.png}
\caption{一幅由 Emanuel Bowen 于 1747 年制作的真实比例太阳系海报。当时，天王星、海王星以及小行星带都尚未被发现。} \label{fig_Planet_18}
\end{figure}
随着科学革命的兴起以及哥白尼、伽利略和开普勒的日心模型的确立，“行星”一词的使用发生了改变：其含义不再是相对于恒星背景在天空中移动的天体，而是指直接（一级行星）或间接（二级行星或卫星行星）绕太阳运行的天体。由此，地球被加入到行星的名单之中 [181]，而太阳被移除。哥白尼体系中的一级行星数量一直保持不变，直到 1781 年威廉·赫歇尔发现天王星 [182]。

当 17 世纪发现木星的四颗卫星（伽利略卫星）以及土星的五颗卫星时，它们与地球的月球一起，被归为绕一级行星运行的“卫星行星”或“二级行星”；不过在随后的几十年里，它们逐渐被简称为“卫星”。直到 20 世纪 20 年代之前，科学家普遍仍将行星的卫星也视为行星，不过这种用法在非科学界并不常见 [154]。

在 19 世纪的第一个十年中，又有四个新的“行星”被发现：谷神星（1801 年）、智神星（1802 年）、婚神星（1804 年）和灶神星（1807 年）。但很快人们意识到，它们与先前已知的行星有明显差异：它们共享同一区域——火星与木星之间的小行星带——并具有部分重叠的轨道。这个区域原本只预期存在一颗行星，而且这些天体的体积远小于其他所有行星；甚至有人怀疑它们可能是某颗更大行星破碎后的碎片。赫歇尔称它们为“小行星”（意为“类星体”），因为即便在最大型的望远镜中，它们也呈现出类似恒星的外观，没有可分辨的盘状结构 [153][183]。

这一局面持续稳定了四十年，但在 1840 年代，又有数颗新的小行星被发现（1845 年的 Astraea；1847 年的 Hebe、Iris 和 Flora；1848 年的 Metis；以及 1849 年的 Hygiea）。新的“行星”几乎每年都有发现；因此，天文学家开始将这些小行星（又称小行星、次级行星）与主要行星分开编目，并为它们分配编号，而不是抽象的行星符号 [153]，尽管它们仍然被视为小型行星 [184]。

海王星于 1846 年被发现，其位置得以预测，是由于它对天王星的引力影响。然而，由于水星的轨道似乎也出现了类似的扰动，19 世纪末人们曾经认为太阳附近可能还存在另一颗行星。然而，水星轨道与牛顿引力预测之间的差异最终由爱因斯坦的广义相对论解释，而不是由另一颗行星解释 [185][186]。

冥王星于 1930 年被发现。由于最初的观测认为它比地球大 [187]，该天体立即被接受为第九颗主要行星。后来更精确的观测发现该天体实际上小得多：1936 年，Ray Lyttleton 提出冥王星可能是海王星的逃逸卫星 [188]；而 Fred Whipple 在 1964 年提出冥王星可能是一颗彗星 [189]。1978 年其大型卫星卡戎（Charon）的发现显示冥王星的质量仅为地球的 0.2\% [190]。由于这仍然远大于任何已知小行星，并且当时尚未发现其他海王星外天体，冥王星保留了行星地位，直到 2006 年才正式失去该地位 [191][192]。

20 世纪 50 年代，Gerard Kuiper 发表了关于小行星起源的论文。他认识到小行星通常并非如此前认为的那样呈球形，并指出小行星族群是碰撞的残余物。因此，他将最大的小行星视为“真正的行星”，而将较小的视为碰撞碎片。从 20 世纪 60 年代开始，“小行星”（asteroid）一词逐渐取代“次级行星”（minor planet），文献中将小行星视为行星的说法也变得罕见，除非讨论地质演化明显的三颗最大天体：谷神星、以及较少被提及的智神星和灶神星 [184]。

20 世纪 60 年代以来，随着空间探测器对太阳系的探索开始，行星科学重新受到重视。此时关于卫星定义的分歧出现：行星科学家开始重新考虑将大型卫星视为行星，而非行星科学的天文学家则通常不这样看待 [154]。（这与前一个世纪的定义并不完全一致，当时所有卫星都被视为次级行星，即使是如土星的 Hyperion 或火星的 Phobos、Deimos 这样非球形的天体也包括在内。）[193][194] 自那以来，八颗主要行星及其行星质量的卫星均已被航天器探测；许多小行星以及矮行星谷神星和冥王星也已被探测。然而，迄今为止，唯一被人类亲自探访的地球以外的行星质量天体仍然是月球。[b]
\subsubsection{界定“行星”一词}
越来越多的天文学家主张将冥王星取消行星地位，因为在 1990 年代和 21 世纪初，在太阳系同一区域（柯伊伯带）发现了许多与其大小相近的天体。事实证明，冥王星只是该区域数千个“小”天体之一 [195]。他们经常引用小行星降格为先例，尽管小行星被降格的原因主要基于其地质物理性质与真正行星的差异，而非因为它们位于小行星带中 [154]。一些较大的海王星外天体，如 Quaoar、Sedna、Eris 和 Haumea，[196] 曾被大众媒体誉为“第十颗行星”。

2005 年 Eris 的发现——其质量比冥王星大 27\%——成为制定行星官方定义的直接动力 [195]。如果将冥王星视为行星，那么逻辑上也必须将 Eris 视为行星。由于行星与非行星的命名程序完全不同，这造成了紧迫的局面：按照当时的规则，在没有界定“行星”的前提下，Eris 是不能被正式命名的 [154]。当时，人们还认为海王星外天体要达到“由自身引力变圆”的尺寸，与巨行星的卫星相近（直径约 400,km）。按照这一标准，柯伊伯带内可能有约 200 个呈圆形的天体，且在更远处可能有数千个 [197][198]。许多天文学家认为，公众不会接受一个导致行星数量剧烈膨胀的定义 [154]。

为解决这一问题，国际天文学联合会（IAU）着手制定“行星”的定义，并于 2006 年 8 月给出了正式版本。在这一定义下，太阳系被认为拥有八颗行星（即水星、金星、地球、火星、木星、土星、天王星和海王星）。那些满足前两条条件但不满足第三条的天体，被归类为矮行星，只要它们不是其他行星的天然卫星。最初，IAU 的一个委员会曾提出一个不包含条件 (c) 的定义方案，这一方案本会将更多的天体纳入行星范畴 [199]。但经过大量讨论，最终通过投票决定，这些天体应被归类为矮行星，而不是行星 [192][200]。

\textbf{对 IAU 定义的批评与替代方案}
\begin{figure}[ht]
\centering
\includegraphics[width=10cm]{./figures/88fe4ac34cbda9b1.png}
\caption{按比例展示的行星质量卫星，并与水星、金星、地球、火星和冥王星进行比较。为了比较，也加入了次行星质量的普罗透斯（Proteus）和奈瑞伊得（Nereid）（其大小与米玛斯相近）。尚未成像的迪斯诺弥娅（Dysnomia）——其大小介于忒堤斯与恩克拉多斯之间——未被展示；无论如何，它很可能并非一个固体天体[90]。} \label{fig_Planet_17}
\end{figure}
IAU 的定义并未被普遍使用或接受。在行星地质学中，天体是依据地球物理特征来被定义为行星的。一个天体在其地幔因自身重量变得具有可塑性的大约质量处，可能会获得一种动力学（行星性的）地质活动。这将导致其进入静水平衡状态，并获得一种稳定的圆形外观，而这种圆形性被地球物理定义视为行星性的标志。例如：[201]

一个具备亚恒星质量、从未发生过核聚变、并且因静水平衡而具有足够引力而呈圆形的天体，不论其轨道参量为何[202]。

在太阳系中，这一质量通常小于清除其轨道所需的质量；因此，在地球物理学定义下被视为“行星”的一些天体（如谷神星和冥王星），在 IAU 定义下并不被视为行星[4]。（在实践中，对静水平衡的要求普遍被放宽为自引力导致的圆整与压实；例如水星实际上并未处于静水平衡状态[203]，但仍然普遍被视为行星[204]。）此类定义的拥护者常常主张，位置不应重要，行星性应由天体的内在性质来界定[4]。矮行星最初被提议作为“小型行星”的一个类别（区别于“类行星体”这样的次行星天体），并且行星地质学家至今仍将它们视为行星，尽管 IAU 定义并非如此[32]。

即便只考虑已知天体，矮行星的数量仍不确定。2019 年，Grundy 等人基于一些中等大小跨海王星天体（TNO）密度较低的事实，提出跨海王星天体达到静水平衡所需的临界尺寸实际上远大于巨行星冰卫星的情形，大约需要直径 900–1000 公里[32]。对于小行星带中的谷神星[205]，以及八个可能跨越此阈值的跨海王星天体——奥库斯（Orcus）、冥王星（Pluto）、妊神星（Haumea）、夸奥尔（Quaoar）、鸟神星（Makemake）、共工（Gonggong）、阋神星（Eris）与赛德娜（Sedna）——普遍存在共识[206][33]。

行星地质学家可能会将十九个已知的行星质量卫星都视为“卫星行星”，包括地球的月球和冥王星的卡戎，就像早期近代天文学家所做的一样[4][207]。一些人甚至更进一步，将如下天体视为行星：尽管今天它们并不是很圆，但具有相对较大、地质演化过的天体，如智神星（Pallas）和灶神星（Vesta）[4]；早期完全被撞击破坏、后又重新聚集并形成圆整外形的天体，如灶神星或许未完全圆整，而节神星（Hygiea）则可能已重新聚并成圆形[208][209][110]；甚至包括所有直径至少达到土卫一（Mimas，最小的行星质量卫星）大小的天体。（这甚至可能包括那些并不圆但尺寸大于土卫一的天体，例如海卫八 Proteus）[4]。

天文学家 Jean-Luc Margot 提出了一项数学判据，可以根据天体质量、轨道半长轴以及母星质量，确定一个天体是否能够在母星生命周期内清除其轨道[210]。该公式产生一个被称为 π 的值，当该值大于 1 时，该天体可视为行星[c]。已知的八大行星以及所有已知系外行星的 π 值都大于 100，而谷神星、冥王星与阋神星的 π 值为 0.1 或更小。π 值达到 1 或更高的天体预计也能大致呈球形，因此对于类太阳恒星而言，满足轨道清除条件的天体通常也会满足圆整条件[211]——尽管在极低质量恒星周围可能并非如此[212]。2024 年，Margot 和合作研究者提出了修订版本，采用统一的轨道清除时间尺度——100 亿年（太阳主序阶段的近似寿命）或 138 亿年（宇宙年龄）——以适用于环绕褐矮星运行的行星[212]。
\subsubsection{系外行星}
在发现系外行星之前，人们就已经对某些天体是否应被视为行星存在特定争议，例如：若它属于某个独立的族群（如某条小天体带），是否仍算行星；又或者它是否足够巨大到能通过氘的热核聚变产生能量[195]。更加复杂的是，质量过小、无法通过氘聚变产生能量的天体，也可能像恒星和棕矮星一样，经由气体云塌缩形成，质量甚至可以低至木星质量[213]；因此，人们对“是否应将天体的形成方式作为行星定义的依据”也存在分歧[195]。

1992 年，天文学家 Aleksander Wolszczan 和 Dale Frail 宣布发现环绕脉冲星 PSR B1257+12 的行星[40]。这一般被视为首个确凿无疑的太阳系外行星系统的发现。随后在 1995 年 10 月 6 日，日内瓦天文台的 Michel Mayor 与 Didier Queloz 公布了首个绕普通主序星运行的系外行星（51 Pegasi）的明确发现[214]。

系外行星的发现又带来了另一项关于行星定义的模糊性：行星与恒星之间的界线在哪里。许多已知的系外行星质量远大于木星，接近被称为棕矮星的恒星类天体。棕矮星通常被视为恒星，因为理论上它们能够进行氘聚变。在银河系中，质量大于约 75 个木星质量的天体可以进行氢聚变，而质量达到 13 个木星质量的天体即可进行氘聚变[215]。然而氘的丰度极低，不足银河系氢含量的 0.0026\%，并且大多数棕矮星在被发现前就早已停止氘聚变，使它们在观测上与超大质量行星几乎无法区分[215]。

\textbf{IAU 对系外行星的工作定义}

2006 年 IAU 的行星定义在应用于系外行星时存在一些困难，因为其措辞明确针对太阳系，而“达到流体静力平衡”以及“清空轨道邻域”这两项判据目前在系外行星上均无法直接观测[1]。随着对系外行星认识的不断增加，该定义在 2018 年进行了重新评估和更新[216]。当前的官方工作定义如下[99]：

\begin{enumerate}
\item 真质量低于氘热核聚变的限制质量（以太阳金属丰度计算，目前为约 13 个木星质量），并且围绕恒星、棕矮星或恒星残骸运行，且与中心天体的质量比低于 L4/L5 不稳定性阈值（即 $M/M_{\text{central}} < 2/(25+\sqrt{621})$）的天体，无论其形成方式如何，均被定义为“行星”。系外行星所需的最低质量/尺寸要求应与太阳系保持一致。
\item 真质量高于氘聚变限制质量的亚恒星天体，无论其形成方式或位置如何，均被定义为“棕矮星”。
\item 在年轻星团中发现的、质量低于氘聚变限制质量的自由漂浮天体不属于“行星”，而被称为“次棕矮星”（或任何更适当的名称）[99]。
\end{enumerate}
IAU 指出，随着知识的改进，这一定义预计将继续演化[99]。一篇 2022 年的综述文章讨论了该定义的历史与合理性，并建议删除第三条中“在年轻星团中”这一措辞，因为此类自由漂浮天体已在其他环境中被发现；同时，建议将“次棕矮星”一词改为更现代、更常用的“自由漂浮行星质量天体”。术语“行星质量天体”也被用于一些具有模糊性的系外行星情境，例如那些质量符合行星范围但却是自由漂浮的，或是围绕棕矮星而非恒星运行的天体[216]。一些自由漂浮的行星质量天体有时仍被称为行星，尤其是“流浪行星”（rogue planets）[217]。

13 个木星质量的氘聚变界线并未得到普遍接受。该界限以下的天体有时也能燃烧氘，而氘的燃烧量取决于天体的内部成分[218][219]。此外，氘的丰度极低，因此氘聚变阶段持续时间很短；与恒星中的氢聚变不同，氘聚变对天体未来演化并无显著影响[58]。质量–半径（或密度）关系在该界限处也没有任何特殊变化，表明棕矮星与较轻的类木行星在物理性质与内部结构上具有连续性，因此从物理角度讲棕矮星被视为行星可能更加自然[58][54]。

因此，许多系外行星目录收录了质量超过 13 个木星质量的天体，有时甚至上限可高达 60 个木星质量[220][100][101][221]。（真正能够进行氢聚变、成为红矮星的质量下限约为 80 个木星质量[58]。）类比主序星的情况，有观点认为“行星”这一术语也应采用类似的包容性定义：主序星在其覆盖的近两个数量级的质量范围内，在结构、温度、大气、光谱特征与可能的形成机制上都有巨大差异，但它们仍被视为同一类天体，因为它们都处于流体静力平衡并进行核聚变[58]。
\subsection{神话与命名}
\subsubsection{古典行星}
英语中太阳系（除 Earth 之外）各行星的名称，源自巴比伦、希腊与罗马在古代逐次发展出的命名传统。将神祇之名赋予行星的做法几乎可以确定是古希腊人从巴比伦人处借鉴而来，其后又被罗马人承袭。巴比伦人将金星以苏美尔的爱情女神命名，其阿卡德语名为 Ishtar；将火星命名为战神 Nergal；将水星命名为智慧之神 Nabu；将木星命名为主神 Marduk；将土星命名为农业之神 Ninurta[222]。希腊人与巴比伦人的命名规则之间存在大量对应关系，因此不可能是各自独立产生的[164]。

鉴于两者神话体系差异显著，这些对应并不完全一致。例如，巴比伦的 Nergal 是战神，因此希腊人将其认同为 Ares。然而，与 Ares 不同，Nergal 也同时是瘟疫之神与冥界统治者[223][224][225]。

在古希腊，两个最明亮的天体——太阳与月亮——分别称为 Helios 与 Selene，皆为古老的泰坦神；最缓慢的行星土星被称为 Phainon，“发光者”；其次是木星，称为 Phaethon，“明亮者”；红色的火星被称为 Pyroeis，“炽焰者”；最明亮的金星称为 Phosphoros，“光之带来者”；最后快速移动的水星称为 Stilbon，“闪耀者”。希腊人将每一颗行星都关联到其神系中的神祇——奥林匹斯诸神或更早的泰坦神[164]：

\begin{itemize}
\item Helios 与 Selene 同时是天体名与神祇名，两者均为泰坦（后被奥林匹斯神 Apollo 与 Artemis 所取代）；
\item Phainon（土星）属于 Cronus（克洛诺斯），这位泰坦是奥林匹斯诸神之父，与丰收相关；
\item Phaethon（木星）属于 Zeus（宙斯），Cronus 之子、推翻其父的王；
\item Pyroeis（火星）属于 Ares（阿瑞斯），宙斯之子、战神；
\item Phosphoros（金星）由 Aphrodite（阿佛洛狄忒）统治，即爱神；
\item Stilbon（水星）因其快速运行而属于 Hermes（赫尔墨斯），即众神使者、学习与机智之神。
\end{itemize}
\begin{figure}[ht]
\centering
\includegraphics[width=8cm]{./figures/2bc830c84ce6a1d1.png}
\caption{} \label{fig_Planet_19}
\end{figure}
尽管现代希腊人仍沿用他们对行星的古代命名，但其他欧洲语言——由于罗马帝国及其后天主教会的影响——采用的却是罗马（拉丁语）的名称，而不是希腊的名称。罗马人与希腊人一样继承了印欧语系原始神话，并以不同名字共享同一套神祇体系；然而，罗马人并不像希腊那样拥有由诗歌文化孕育出的丰富神话叙事传统。在罗马共和国的后期，罗马作家大量借用希腊神话故事并将其应用于罗马自身的神祇体系，最终两者几乎难以区分[226]。当罗马人研究希腊天文学时，他们用自己神祇的名字来重新命名行星：Mercurius（对应 Hermes）、Venus（Aphrodite）、Mars（Ares）、Iuppiter（Zeus）与 Saturnus（Cronus）。然而，对于某颗行星应对应哪位神祇，人们并未达成广泛一致。根据老普林尼的记载，虽然 Phainon 与 Phaethon 分别对应 Saturn 与 Jupiter 已得到一致认可，但 Pyroeis 也被关联到半神 Hercules；Stilbon 也被关联到音乐、医治与预言之神 Apollo；Phosphoros 甚至同时被关联到重要女神 Juno 与 Isis[227]。部分罗马人——遵循一种可能起源于美索不达米亚、但在希腊化时期的埃及得到发展的信念——认为七位以行星命名的神祇轮流按小时守护地上事务。轮值顺序为 Saturn、Jupiter、Mars、Sun、Venus、Mercury、Moon（从最远行星到最近）[228]。因此：第一天由 Saturn 开始（第 1 小时），第二天由 Sun 开始（第 25 小时），第三天由 Moon 开始（第 49 小时），接着依次为 Mars、Mercury、Jupiter、Venus。由于每一天以开启该日的神祇命名，这就形成了罗马历法中一周的顺序[229]。在英语中，Saturday、Sunday 与 Monday 是这些罗马名称的直接翻译。而其他几天则以盎格鲁-撒克逊诸神命名：Tīw（Tuesday）、Wōden（Wednesday）、Þunor（Thursday）、Frīġ（Friday），分别被认为与 Mars、Mercury、Jupiter 与 Venus 相似或对应[230]。

地球在英语中的名称并非源自希腊—罗马神话。由于直到 17 世纪地球才被普遍接受为一颗行星[181]，因此并没有把它命名为某位神祇的传统。（在英语中，太阳与月亮也属于这种情况，尽管它们如今一般不再被视为行星。）“Earth”这一名称源自古英语 eorþe，该词既指“土地”“泥土”，也指整个世界[231]。与其他日耳曼语言中的同源词一样，它最终源于原始日耳曼语的 erþō，其痕迹仍可在现代英语 earth、德语 Erde、荷兰语 aarde 与斯堪的纳维亚语 jord 中看到。许多罗曼语族语言则保留了古罗马词汇 terra（及其变体），其含义是“干地”而非“海洋”[232]。非罗曼语族语言则使用其自身的本源词。希腊语仍使用其原始名称 Γή（Ge）[233]。

非欧洲文化使用其他的行星命名体系。印度采用基于九曜（Navagraha）的体系，包括七颗传统行星及上升、下降的月球交点 Rahu 与 Ketu。各行星名称如下：Surya（太阳）、Chandra（月亮）、Budha（对应 Mercury）、Shukra（“明亮”之意，对应 Venus）、Mangala（战神，对应 Mars）、Bṛhaspati（众神顾问，对应 Jupiter）、Shani（象征时间，对应 Saturn）[234]。

大多数行星的波斯本土名称基于美索不达米亚诸神与伊朗诸神的对应关系，与希腊及拉丁传统类似。Mercury 为 Tir（波斯语：تیر），对应西伊朗文书吏守护神 Tīriya（类似 Nabu）；Venus 为 Nāhid（ناهید，对应 Anahita）；Mars 为 Bahrām（بهرام，对应 Verethragna）；Jupiter 为 Hormoz（هرمز，对应 Ahura Mazda）。Saturn 的波斯名称 Keyvān（کیوان）借自阿卡德语 kajamānu，意为“永恒、稳定”[235]。

中国与受中国文化影响的东亚国家（包括日本、韩国、越南）采用五行命名体系：水（Mercury 水星）、金（Venus 金星）、火（Mars 火星）、木（Jupiter 木星）、土（Saturn 土星）[229]。

在传统希伯来天文学中，七颗古典行星（大部分）具有描述性名称——太阳为 חמה Ḥammah（“炎热者”），月亮为 לבנה Levanah（“洁白者”），Venus 为 כוכב נוגה Kokhav Nogah（“明亮之星”），Mercury 为 כוכב Kokhav（“行星”，因其缺乏显著特征），Mars 为 מאדים Ma'adim（“红者”），Saturn 为 שבתאי Shabbatai（“休息者”，指其相较其他可见行星运动缓慢）[236]。唯一的例外是 Jupiter，被称为 צדק Tzedeq（“正义”）[236]。这些名称首次见于《巴比伦塔木德》，并非希伯来语中对行星的原始称呼。公元 377 年，萨拉米斯的 Epiphanius 记录了另一套名称，这些名称似乎具有异教或迦南语背景；其后因宗教原因被替换，但它们很可能是历史上闪米特族对行星的真实命名，甚至可能追溯到更早的巴比伦天文学[236]。

阿拉伯语行星名称的词源理解得不那么充分。学界普遍同意的只有以下几项：Venus（Arabic: الزهرة, az-Zuhara, “明亮者”[237]）、Earth（الأرض, al-ʾArḍ，与希伯来语 eretz 同源）、Saturn（زُحَل, Zuḥal, “退离者”[238]）。至于 Mercury（عُطَارِد, ʿUṭārid）、Mars（اَلْمِرِّيخ, al-Mirrīkh）与 Jupiter（المشتري, al-Muštarī）的词源，学界有多种不同假说，但尚无共识[239][240][241][242]。
\subsubsection{现代时期的发现}
当 18 与 19 世纪又有新的行星被发现时，天王星以希腊神祇命名，而海王星以罗马神祇（即希腊波塞冬的对应者）命名。小行星最初也采用神话命名——谷神星（Ceres）、婚神星（Juno）、灶神星（Vesta）都是主要的罗马女神，而智神星（Pallas）则是希腊主神雅典娜的称号之一——但随着新小行星不断被发现，人们先开始转向较次要的女神命名，直到第 20 号小行星 Massalia（1852 年）为止，神话命名限制才正式被舍弃[243]。冥王星（以希腊冥界之神命名）获得一个古典名称，是因为在其被发现时，它被视为一颗主要行星。

中文、韩文、日文中的天王星（“天王星”、韓：천왕성、日：天王星）、海王星（“海王星”）、冥王星（“冥王星”）都是基于希腊—罗马神话中相关神祇职能的意译[244][245][d]。19 世纪，伟烈亚力（Alexander Wylie）与李善兰共同将最初发现的 117 颗小行星译为中文，其中许多名称沿用至今，如：谷神星（穀神星）、智神星、婚神星、灶神星与健神星（Hygiea）[247]。此类翻译后来扩展到了部分 21 世纪发现的矮行星，例如：妊神星（Haumea）、鳥神星（Makemake）与鬩神星（Eris）。然而，除了这些较有知名度的对象，大多数译名在中文天文学文献之外相当少见[244]。

2009 年，人们为天王星与海王星正式选定了希伯来语名称——天王星为 אורון Oron,（“小光”），海王星为 רהב Rahab（圣经中的海怪）[248]；在此之前，希伯来语中直接借用 “Uranus” 与 “Neptune”[249]。

随着海王星外天体大量被发现，与轨道相关的命名规则开始建立：与海王星存在 2:3 共振的天体（冥族小天体，plutinos）通常以冥界神话命名，而其他轨道族群多以创世神话命名。大多数海王星外天体的名称来自世界各文化的神祇（如 Quaoar 以通格瓦族神祇命名）。少数例外延续希腊—罗马传统，尤其是 Eris，因为它曾一度被视为“第十行星”[250][251]。

行星的卫星（包括那些达行星质量的卫星）通常以与母行星相关的神话体系命名。木星的行星质量卫星以前称为“宙斯的情人或性伴侣”；土星的卫星以克罗诺斯的兄弟姐妹——泰坦众神——命名；天王星的卫星以莎士比亚与蒲柏作品中的角色命名（最初限定于“精灵神话”角色[252]，但自米兰达命名后不再坚持此规则）。海王星的行星质量卫星 Triton 以其子命名；冥王星的行星质量卫星 Charon 则以冥界摆渡者命名——他负责运送亡灵进入冥王的领域[253]。
\subsubsection{系外行星}
系外行星通常以其母恒星名称加上发现顺序的字母表示，例如 Proxima Centauri b。（字母从 b 开始，a 被视为指代恒星本身。）[citation needed]
\subsubsection{符号}
\begin{figure}[ht]
\centering
\includegraphics[width=14.25cm]{./figures/8435d2047d6e74d6.png}
\caption{} \label{fig_Planet_20}
\end{figure}
水星、金星、木星、土星以及可能的火星的书写符号可以追溯到希腊晚期纸草文献中的形体[254]。木星与土星的符号被认为是其希腊名称的合字，而水星的符号是一种风格化的商神杖（caduceus）[254]。

根据 Annie Scott Dill Maunder 的研究，行星符号的前身曾在艺术中用于代表与古典行星相关的神祇。Bianchini 的天球仪（由 Francesco Bianchini 于 18 世纪发现，但制作于 2 世纪）[255] 展示了带有早期行星符号的希腊行星神拟人像：水星持商神杖；金星的项链上连着另一条绳饰；火星持长矛；木星持权杖；土星持镰刀；太阳呈放射状光环；月亮戴有新月饰的头饰[256]。带有十字标记的现代行星符号大约在 16 世纪首次出现。按照 Maunder 的解释，这些十字似乎是“试图给旧异教诸神的符号增添某种基督教意味”[256]。地球本身未被视为古典行星；其符号则源自日心说出现前，用以表示“四方世界”的象征[257]。

当更多绕日运行的行星被发现后，人们又为它们创造新的符号。天王星最常见的天文符号 ⛢[258] 由 Johann Gottfried Köhler 发明，旨在象征当时新发现的金属——铂金[259][260]。另一种符号 ♅ 则由 Jérôme Lalande 创造，表示顶部带有“H”（纪念发现者 Herschel）的球体[261]。如今，⛢ 通常用于天文学，而 ♅ 多用于占星学，但两者偶尔也会交叉出现[258]。

最初被发现的几颗小行星当时被视为行星，因此也被赋予象征符号，例如：谷神星的镰刀（⚳）、智神星的长矛（⚴）、婚神星的权杖（⚵）、灶神星的炉灶标志（⚶）。但随着数量快速增加，这种做法被编号系统取代。（Massalia 是第一颗不以神话命名的小行星，也因此成为第一颗未被赋予象征符号的小行星。）前四颗小行星（谷神星至灶神星）的符号保留使用的时间比其余小行星长得多[153]；在现代，NASA 仍偶尔使用谷神星符号——谷神星是唯一也是矮行星的小行星[262]。

海王星的符号（♆）代表该神明的三叉戟[260]。冥王星的天文符号为 P-L 合字（♇）[263]，但自 IAU 将冥王星重新分类后，此符号日趋少用[262]。自冥王星被重新归类为矮行星后，NASA 使用的是冥王星的传统占星符号（⯓），即顶端为天体圆形、下方为冥王之双叉杖的图案[262]。
\begin{figure}[ht]
\centering
\includegraphics[width=14.25cm]{./figures/c2c6c35667c42ca9.png}
\caption{} \label{fig_Planet_21}
\end{figure}
国际天文学联合会（IAU）不鼓励在现代期刊文章中使用行星符号，而是主张对主要行星采用单字母缩写，或（为区分水星与火星）使用双字母缩写。[264]
然而，太阳与地球的符号仍然常见，因为诸如太阳质量、地球质量等单位在天文学中应用广泛。[264]

其他行星符号如今主要出现在占星学中。占星家重新启用了前几颗小行星的古老天文符号，并继续为更多天体创造新符号。[262]
这当中包括 21 世纪发现的矮行星的占星符号——这些天体的发现时代已经晚于行星符号在天文学中普遍弃用的时期，因此天文学家并未为它们分配新符号。[262]
许多占星符号已被纳入 Unicode，其中部分较新的符号（如妊神星 Haumea、鸟神星 Makemake、鬩神星 Eris 的符号）后来也被 NASA 用于天文学场景。[262]

Eris 的符号源自 *Discordianism*（不和谐主义宗教）中对女神 Eris 的传统符号。[262]
其他矮行星的符号（除 Haumea 外）大多是其母文化文字体系中的首字母组合，同时兼具与该神祇或文化意象相关的象征意义，例如：
Makemake 的符号代表神祇的面容、共工（Gonggong）的符号代表其蛇尾。[262][265]

Moskowitz 也为行星质量级别的卫星设计了符号；多数符号由首字母与其母行星符号的某个元素组合而成。唯一的例外是冥卫一卡戎（Charon）的符号：它将冥王星双叉杖符号上的高位圆环与一弯新月结合，象征了卡戎作为卫星的身份，同时也隐喻神话中卡戎划船渡过冥河斯堤克斯（Styx）。[266]
\subsection{参见}
\begin{itemize}
\item \textbf{Double planet} —— 双行星；指两个具行星质量的天体，其公共轨道轴心位于两者之外的双体系统。
\item \textbf{List of landings on extraterrestrial bodies} —— 各类航天器登陆地外天体的列表。
\item \textbf{Lists of planets} —— 按不同属性分类的行星列表集合。
\item \textbf{Mesoplanet} —— 提议用语，指质量介于水星与谷神星之间的行星级天体。
\item \textbf{Planetary habitability} —— 行星宜居性；指行星适合生命存在的已知程度。
\item \textbf{Planetary mnemonic} —— 用于记忆太阳系行星顺序的助记短语。
\item \textbf{Theoretical planetology} —— 行星的理论模型研究。
\end{itemize}
\subsection{注释}\\
a.这里“类地大小”指半径为地球半径的 1–2 倍，“宜居带”指接收恒星辐射量为地球辐射量 0.25 至 4 倍的区域（对应于太阳的 0.5–2 AU）。目前没有关于太阳这类 G 型恒星的统计数据；该统计基于 K 型恒星的数据外推。[51][52]\\
b.参见：Solar System exploration 的时间线。\\
c.Margot 的参数[211]不应与著名的数学常数 $\pi\approx 3.14159265\ldots$ 混淆。\\
d.在韩语中，这些行星名称更常以韩文书写，而非汉字，例如冥王星写作 명왕성。在越南语中，比起直接采用汉越音，更常使用意译的词语，例如水星为 sao Thuỷ（而非汉越音的 Thuỷ tinh）。冥王星不是 sao Minh Vương，而是 sao Diêm Vương（“阎王星”，对应阎魔天）。[246]
\subsection{参考文献}
\begin{enumerate}
\item Lecavelier des Etangs, A.; Lissauer, Jack J.（2022 年 6 月 1 日），“IAU 对系外行星的工作定义”，New Astronomy Reviews，94：101641。arXiv:2203.09520。Bibcode:2022NewAR..9401641L。doi:10.1016/j.newar.2022.101641。ISSN 1387-6473。S2CID 247065421。于 2022 年 5 月 13 日存档。2022 年 5 月 13 日查阅。
\item “IAU 2006 年大会：IAU 决议投票结果”。国际天文学联合会（IAU）。2006。2004 年 4 月 29 日存档。2009 年 12 月 30 日查阅。
\item “国际天文学联合会系外行星工作组（WGESP）”。IAU。2001。2006 年 9 月 16 日存档。2008 年 8 月 23 日查阅。
\item Lakdawalla, Emily（2020 年 4 月 21 日），“什么是行星？”，The Planetary Society。2022 年 1 月 22 日存档。2022 年 4 月 3 日查阅。
\item Grossman, Lisa（2021 年 8 月 24 日），“行星定义仍然是敏感话题——尤其是对冥王星支持者而言”。Science News。2022 年 7 月 10 日存档。2022 年 7 月 10 日查阅。
\item Wetherill, G. W.（1980），“类地行星的形成”，Annual Review of Astronomy and Astrophysics，18（1）：77–113。Bibcode:1980ARA&A..18...77W。doi:10.1146/annurev.aa.18.090180.000453。ISSN 0066-4146。
\item D'Angelo, G.; Bodenheimer, P.（2013），“嵌入原行星盘中的年轻行星包层的三维辐射流体动力学计算”，The Astrophysical Journal，778（1）：77（29 页）。arXiv:1310.2211。Bibcode:2013ApJ...778...77D。doi:10.1088/0004-637X/778/1/77。
\item Inaba, S.; Ikoma, M.（2003），“拥有大气层的原行星的增强碰撞增长”，Astronomy and Astrophysics，410（2）：711–723。Bibcode:2003A&A...410..711I。doi:10.1051/0004-6361:20031248。
\item D'Angelo, G.; Weidenschilling, S. J.; Lissauer, J. J.; Bodenheimer, P.（2014），“木星的增长：低质量包层增强核心吸积作用”，Icarus，241：298–312。arXiv:1405.7305。Bibcode:2014Icar..241..298D。doi:10.1016/j.icarus.2014.06.029。
\item Lissauer, J. J.; Hubickyj, O.; D'Angelo, G.; Bodenheimer, P.（2009），“结合热力学与流体力学约束的木星增长模型”，Icarus，199（2）：338–350。arXiv:0810.5186。Bibcode:2009Icar..199..338L。doi:10.1016/j.icarus.2008.10.004。
\item D'Angelo, G.; Durisen, R. H.; Lissauer, J. J.（2011），“巨行星的形成”，载于 Seager, S.（编），Exoplanets，亚利桑那大学出版社（图森），页 319–346。arXiv:1006.5486。Bibcode:2010exop.book..319D。2015 年 6 月 30 日存档。2016 年 5 月 1 日查阅。
\item Chambers, J.（2011），“类地行星形成”，载于 Seager, S.（编），Exoplanets，亚利桑那大学出版社，图森，页 297–317。Bibcode:2010exop.book..297C。2015 年 6 月 30 日存档。2016 年 5 月 1 日查阅。
\item Canup, Robin M.; Ward, William R.（2008），Origin of Europa and the Galilean Satellites，亚利桑那大学出版社，第 59 页。arXiv:0812.4995。Bibcode:2009euro.book...59C。ISBN 978-0-8165-2844-8。
\item D'Angelo, G.; Podolak, M.（2015），“木星周围盘中的微行星的捕获与演化”，The Astrophysical Journal，806（1）：29 页。arXiv:1504.04364。Bibcode:2015ApJ...806..203D。doi:10.1088/0004-637X/806/2/203。
\item Agnor, C. B.; Hamilton, D. P.（2006），“海王星在双行星引力遭遇中的特里同捕获”（PDF），Nature，441（7090）：192–4。Bibcode:2006Natur.441..192A。doi:10.1038/nature04792。PMID 16688170。S2CID 4420518。2016 年 10 月 14 日 PDF 存档。2022 年 5 月 1 日查阅。
\item Taylor, G. Jeffrey（1998 年 12 月 31 日），“地球与月球的起源”，Planetary Science Research Discoveries，夏威夷地球物理与行星学研究所。2010 年 6 月 10 日存档。2010 年 4 月 7 日查阅。
\item Stern, S. A.; Bagenal, F.; Ennico, K.; Gladstone, G. R.; 等（2015 年 10 月 16 日），“冥王星系统：新视野号探测的初步结果”，Science，350（6258）：aad1815。arXiv:1510.07704。Bibcode:2015Sci...350.1815S。doi:10.1126/science.aad1815。PMID 26472913。S2CID 1220226。
\item Dutkevitch, Diane（1995），年轻恒星周围类地行星区域尘埃的演化（博士论文），马萨诸塞大学阿默斯特分校。Bibcode:1995PhDT..........D。2007 年 11 月 25 日存档。2008 年 8 月 23 日查阅。
\item Matsuyama, I.; Johnstone, D.; Murray, N.（2005），“由中心源光致蒸发抑制行星迁移”，The Astrophysical Journal，585（2）：L143–L146。arXiv:astro-ph/0302042。Bibcode:2003ApJ...585L.143M。doi:10.1086/374406。
\item Kenyon, Scott J.; Bromley, Benjamin C.（2006），“类地行星形成 I：从寡头增长到混沌增长的过渡”，Astronomical Journal，131（3）：1837–1850。arXiv:astro-ph/0503568。Bibcode:2006AJ....131.1837K。doi:10.1086/499807。
\item Martin, R. G.; Livio, M.（2013 年 1 月 1 日），“关于小行星带的形成与演化及其对生命的潜在意义”，Monthly Notices of the Royal Astronomical Society: Letters，428（1）：L11–L15。arXiv:1211.0023。doi:10.1093/mnrasl/sls003。ISSN 1745-3925。
\item Peale, S. J.（1999 年 9 月），“天然卫星的起源与演化”，Annual Review of Astronomy and Astrophysics，37（1）：533–602。Bibcode:1999ARA&A..37..533P。doi:10.1146/annurev.astro.37.1.533。ISSN 0066-4146。2022 年 5 月 13 日存档。2022 年 5 月 13 日查阅。
\item Ida, Shigeru; Nakagawa, Yoshitsugu; Nakazawa, Kiyoshi（1987），“雷利–泰勒不稳定导致的地核形成”，Icarus，69（2）：239–248。Bibcode:1987Icar...69..239I。doi:10.1016/0019-1035(87)90103-5。
\item Kasting, James F.（1993），“地球早期大气”，Science，259（5097）：920–926。Bibcode:1993Sci...259..920K。doi:10.1126/science.11536547。PMID 11536547。S2CID 21134564。
\item Chuang, F.（2012 年 6 月 6 日），“常见问题——大气层”，Planetary Science Institute。2022 年 3 月 23 日存档。2022 年 5 月 13 日查阅。
\item Fischer, Debra A.; Valenti, Jeff（2005），“行星—金属量关联”，The Astrophysical Journal，622（2）：1102。Bibcode:2005ApJ...622.1102F。doi:10.1086/428383。
\item Wang, Ji; Fischer, Debra A.（2013），“揭示太阳型恒星周围不同大小行星普适的金属量关联”，The Astronomical Journal，149（1）：14。arXiv:1310.7830。Bibcode:2015AJ....149...14W。doi:10.1088/0004-6256/149/1/14。S2CID 118415186。
\item Harrison, Edward Robert（2000），Cosmology: The Science of the Universe，剑桥大学出版社，第 114 页。ISBN 978-0-521-66148-5。2023 年 12 月 14 日存档。2022 年 5 月 13 日查阅。
\item “行星的物理参数”，Solar System Dynamics，喷气推进实验室（JPL）。2022 年 10 月 4 日存档。2022 年 7 月 11 日查阅。
\item Lewis, John S.（2004），Physics and Chemistry of the Solar System（第 2 版），Academic Press，第 59 页。ISBN 978-0-12-446744-6。
\item Marley, Mark（2019 年 4 月 2 日），“并非冰之心”，planetary.org，The Planetary Society。2019 年 8 月 12 日存档。2022 年 5 月 5 日查阅。
\item Grundy, W. M.; Noll, K. S.; Buie, M. W.; Benecchi, S. D.; 等（2018 年 12 月），“跨海王星双星 Gǃkúnǁʼhòmdímà（(229762) 2007 UK126）的互绕轨道、质量与密度”（PDF），Icarus，334：30。Bibcode:2019Icar..334...30G。doi:10.1016/j.icarus.2018.12.037。S2CID 126574999。2019 年 4 月 7 日存档。
\item Emery, J. P.; Wong, I.; Brunetto, R.; Cook, J. C.; Pinilla-Alonso, N.; 等（2024），“三颗矮行星的故事：JWST 光谱揭示 Sedna、Gonggong 与 Quaoar 上的冰与有机物”，Icarus，414。arXiv:2309.15230。Bibcode:2024Icar..41416017E。doi:10.1016/j.icarus.2024.116017。
\item Brown, Michael E.; Schaller, Emily L.（2007 年 6 月 15 日），“矮行星阋神星的质量”（PDF），Science，316（5831）：1585。Bibcode:2007Sci...316.1585B。doi:10.1126/science.1139415。PMID 17569855。S2CID 21468196。2016 年 3 月 4 日 PDF 存档。2015 年 9 月 27 日查阅。
\item “冥王星有多大？新视野号终结了持续数十年的争论”，NASA（2017 年 8 月 7 日）。2019 年 11 月 9 日存档。2022 年 5 月 5 日查阅。
\item Lewis, John S.（2004），Physics and Chemistry of the Solar System（第 2 版），Academic Press，第 425 页。ISBN 978-0-12-446744-6。
\item “预生成系外行星图表”，exoplanetarchive.ipac.caltech.edu，NASA Exoplanet Archive。2012 年 4 月 30 日存档。2022 年 6 月 24 日查阅。
\item Schneider, J.，“交互式系外行星目录”，The Extrasolar Planets Encyclopedia。2025 年 10 月 30 日查阅。
\item Cassan, Arnaud; Kubas, D.; Beaulieu, J.-P.; Dominik, M.; 等（2012 年 1 月 12 日），“来自微引力观测的银河系恒星平均具有一颗或更多束缚行星”，Nature，481（7380）：167–169。arXiv:1202.0903。Bibcode:2012Natur.481..167C。doi:10.1038/nature10684。PMID 22237108。S2CID 2614136。
\item Wolszczan, A.; Frail, D. A.（1992），“位于毫秒脉冲星 PSR1257+12 周围的行星系统”，Nature，355（6356）：145–147。Bibcode:1992Natur.355..145W。doi:10.1038/355145a0。S2CID 4260368。
\item Wolszczan, Alex（2008），“PSR B1257+12 脉冲星周围的行星”，Extreme Solar Systems，398：3+。Bibcode:2008ASPC..398....3W。2022 年 5 月 13 日存档。2022 年 5 月 13 日查阅。
\item “外面有什么世界？”，加拿大广播公司（CBC），2016 年 8 月 25 日。2016 年 8 月 25 日存档。2017 年 6 月 5 日查阅。
\item Chen, Rick（2018 年 10 月 23 日），“开普勒任务的顶级科学成果”，NASA。2022 年 7 月 11 日存档。2022 年 7 月 11 日查阅。开普勒发现的最常见行星大小是介于地球与海王星之间、太阳系中不存在的类型，我们仍需进一步了解这类行星。
\item Barclay, Thomas; Rowe, Jason F.; Lissauer, Jack J.; Huber, Daniel; 等（2013 年 2 月 28 日），“一颗亚水星大小的系外行星”，Nature，494（7438）：452–454。arXiv:1305.5587。Bibcode:2013Natur.494..452B。doi:10.1038/nature11914。ISSN 0028-0836。PMID 23426260。S2CID 205232792。2022 年 10 月 19 日存档。2022 年 7 月 11 日查阅。
\item Johnson, Michele（2011 年 12 月 20 日），“NASA 发现太阳系外首批地球大小行星”，NASA。2020 年 5 月 16 日存档。2011 年 12 月 20 日查阅。
\item Hand, Eric（2011 年 12 月 20 日），“开普勒发现首批地球大小的系外行星”，Nature。doi:10.1038/nature.2011.9688。S2CID 122575277。
\item Overbye, Dennis（2011 年 12 月 20 日），“两颗地球大小的行星被发现”，纽约时报。2011 年 12 月 20 日存档。2011 年 12 月 21 日查阅。
\item Kopparapu, Ravi Kumar（2013），“对开普勒 M 型矮星宜居带类地行星出现率的修正估计”，The Astrophysical Journal Letters，767（1）：L8。arXiv:1303.2649。Bibcode:2013ApJ...767L...8K。doi:10.1088/2041-8205/767/1/L8。S2CID 119103101。
\item Watson, Traci（2016 年 5 月 10 日），“NASA 的发现使已知行星数量翻倍”，USA Today。2016 年 5 月 10 日存档。2016 年 5 月 10 日查阅。
\item “宜居系外行星目录”，波多黎各阿雷西博大学，行星宜居性实验室。2011 年 10 月 20 日存档。2022 年 7 月 12 日查阅。
\item Sanders, R.（2013 年 11 月 4 日），“天文学家回答关键问题：宜居行星有多常见？”，伯克利新闻中心。2014 年 11 月 7 日存档。2013 年 11 月 7 日查阅。
\item Petigura, E. A.; Howard, A. W.; Marcy, G. W.（2013），“围绕类太阳恒星运行的地球大小行星的普遍性”，PNAS，110（48）：19273–19278。arXiv:1311.6806。Bibcode:2013PNAS..11019273P。doi:10.1073/pnas.1319909110。PMC 3845182。PMID 24191033。
\item Drake, Frank（2003 年 9 月 29 日），“德雷克方程再探”，Astrobiology Magazine。2011 年 6 月 28 日存档。2008 年 8 月 23 日查阅。
\item Chen, Jingjing; Kipping, David（2016），“其他世界的质量与半径的概率预测”，The Astrophysical Journal，834（1）：17。arXiv:1603.08614。Bibcode:2017ApJ...834...17C。doi:10.3847/1538-4357/834/1/17。S2CID 119114880。
\item Mayor, Michel; Bonfils, Xavier; Forveille, Thierry; 等（2009），“HARPS 南天系外行星巡天 XVIII：GJ 581 行星系统中的一颗类地质量行星”（PDF），Astronomy and Astrophysics，507（1）：487–494。arXiv:0906.2780。Bibcode:2009A&A...507..487M。doi:10.1051/0004-6361/200912172。S2CID 2983930。2009 年 5 月 21 日 PDF 存档。
\item “发现新的‘超级地球’”，BBC News，2007 年 4 月 25 日。2012 年 11 月 10 日存档。2007 年 4 月 25 日查阅。
\item von Bloh 等（2007），“Gliese 581 系统中超级地球的宜居性”，Astronomy and Astrophysics，476（3）：1365–1371。arXiv:0705.3758。Bibcode:2007A&A...476.1365V。doi:10.1051/0004-6361:20077939。S2CID 14475537。
\item Hatzes, Artie P.; Rauer, Heike（2015），“基于质量—密度关系的巨行星定义”，The Astrophysical Journal，810（2）：L25。arXiv:1506.05097。Bibcode:2015ApJ...810L..25H。doi:10.1088/2041-8205/810/2/L25。S2CID 119111221。
\item Zhang, Zhoujian; Liu, Michael C.; Claytor, Zachary R.; Best, William M. J.; 等（2021 年 8 月 1 日），“COCONUTS 计划的第二项发现：一颗在 10.9 pc 处年轻场 M 矮星外侧的寒冷宽轨系外行星”，The Astrophysical Journal Letters，916（2）：L11。arXiv:2107.02805。Bibcode:2021ApJ...916L..11Z。doi:10.3847/2041-8213/ac1123。hdl:20.500.11820/4f26e8e5-5d42-4259-bc20-fcb093d664b6。ISSN 2041-8205。S2CID 236464073。
\item “系外行星”，科罗拉多大学 LASP。2019 年 4 月 5 日存档。2022 年 5 月 13 日查阅。

\item Anderson, D. R.; Hellier, C.; Gillon, M.; Triaud, A. H. M. J.; 等（2009），“WASP-17b：一颗极低密度并处于可能的逆行轨道上的行星”，\textit{The Astrophysical Journal}，709（1）：159–167。arXiv:0908.1553。Bibcode:2010ApJ...709..159A。doi:10.1088/0004-637X/709/1/159。S2CID 53628741。

\item Young, Charles Augustus（1902），\textit{Manual of Astronomy: A Text Book}，Ginn & Company，324–327 页。

\item Dvorak, R.; Kurths, J.; Freistetter, F.（2005），\textit{Chaos And Stability in Planetary Systems}，纽约：Springer，第 90 页。ISBN 978-3-540-28208-2。

\item Moorhead, Althea V.; Adams, Fred C.（2008），“巨行星轨道偏心率在环星盘力矩作用下的演化”，\textit{Icarus}，193（2）：475–484。arXiv:0708.0335。Bibcode:2008Icar..193..475M。doi:10.1016/j.icarus.2007.07.009。S2CID 16457143。

\item “行星——柯伊伯带天体”，\textit{The Astrophysics Spectator}，2004 年 12 月 15 日。2021 年 3 月 23 日存档。2008 年 8 月 23 日查阅。

\item Tatum, J. B.（2007），“17. 目视双星”，\textit{Celestial Mechanics}，个人主页。2007 年 7 月 6 日存档。2008 年 2 月 2 日查阅。

\item Trujillo, Chadwick A.; Brown, Michael E.（2002），“古典柯伊伯带中轨道倾角与颜色的相关性”，\textit{Astrophysical Journal}，566（2）：L125。arXiv:astro-ph/0201040。Bibcode:2002ApJ...566L.125T。doi:10.1086/339437。S2CID 11519263。

\item Peter Goldreich（1966 年 11 月），“月球轨道历史”，\textit{Reviews of Geophysics}，4（4）：411–439。Bibcode:1966RvGSP...4..411G。doi:10.1029/RG004i004p00411。

\item Harvey, Samantha（2006 年 5 月 1 日），“天气，到处都是天气？”，NASA。2006 年 8 月 31 日存档。2008 年 8 月 23 日查阅。

\item 行星数据表，NASA。

\item Schorghofer, N.; Mazarico, E.; Platz, T.; Preusker, F.; Schröder, S. E.; Raymond, C. A.; Russell, C. T.（2016 年 7 月 6 日），“矮行星谷神星的永久阴影区”，\textit{Geophysical Research Letters}，43（13）：6783–6789。Bibcode:2016GeoRL..43.6783S。doi:10.1002/2016GL069368。

\item Carry, B.; 等（2009），“（2）智神星的物理性质”，Icarus，205（2）：460–472。arXiv:0912.3626。Bibcode:2010Icar..205..460C。doi:10.1016/j.icarus.2009.08.007。S2CID 119194526。
\item Thomas, P. C.; 等（1997），“灶神星的自转轴、大小与形状（来自 HST 图像）”，\textit{Icarus}，128（1）：88–94。Bibcode:1997Icar..128...88T。doi:10.1006/icar.1997.5736。
\item Winn, Joshua N.; Holman, Matthew J.（2005），“热木星上的倾角潮汐”，\textit{The Astrophysical Journal}，628（2）：L159。arXiv:astro-ph/0506468。Bibcode:2005ApJ...628L.159W。doi:10.1086/432834。S2CID 7051928。
\item Seidelmann, P. Kenneth（编）（1992），Explanatory Supplement to the Astronomical Almanac，University Science Books，第 384 页。
\item Lang, Kenneth R.（2011），The Cambridge Guide to the Solar System（第 2 版），剑桥大学出版社。ISBN 978-1-139-49417-5。2016 年 1 月 1 日存档。
\item Bills, Bruce G.（2005），“木星伽利略卫星的自由与强迫自转轴倾角”，Icarus，175（1）：233–247。Bibcode:2005Icar..175..233B。doi:10.1016/j.icarus.2004.10.028。2020 年 7 月 27 日存档。2023 年 4 月 6 日查阅。
\item Goldstein, R. M.; Carpenter, R. L.（1963），“金星自转：来自雷达测量的周期估计”，Science，139（3558）：910–911。Bibcode:1963Sci...139..910G。doi:10.1126/science.139.3558.910。PMID 17743054。S2CID 21133097。
\item Belton, M. J. S.; Terrile, R. J.（1984），收录于 Bergstralh, J. T.（编），Uranus–Neptune 研讨会（1984 年 2 月 6–8 日，帕萨迪纳）：天王星与海王星的自转特性，327–347 页。Bibcode:1984NASCP2330..327B。
\item Borgia, Michael P.（2006），The Outer Worlds; Uranus, Neptune, Pluto, and Beyond，Springer New York，195–206 页。



\end{enumerate}








